\section{Introduction}
\textcolor{red}{koushik}

Cancer treatment planning is a difficult clinical problem because one
patient may need several different kinds of decisions at the same time. A
clinician may need to decide whether a patient should receive a targeted
molecular therapy, whether radiation therapy should be used, and whether
the patient is likely to survive beyond a short-term clinical window. These
three questions are related, but they are not the same question. Targeted
molecular therapy aims to block specific molecular drivers of cancer
growth~\cite{Min2022}. Radiation therapy uses ionising radiation to damage
tumour cells~\cite{Baskar2012}. Short-term survival prediction is often used
to support risk stratification and care planning in oncology~\cite{Parikh2019,
SideyGibbons2022}. A useful model should therefore learn these outcomes as
separate but connected prediction tasks, instead of forcing them into one
single label.

Gene expression data can help with this problem because it captures which
genes are active in each tumour sample. However, raw gene expression has
thousands of genes, and it can be hard to understand what a model learns
from such a large input. A more interpretable approach is to summarize gene
expression into biological pathways. Gene set enrichment analysis was
introduced to connect gene expression patterns with known biological
processes~\cite{subramanian2005gsea}, and single-sample GSEA makes it
possible to estimate pathway activity for each individual
sample~\cite{barbie2009ssgsea}. In this paper, we use Reactome, a curated
pathway knowledgebase, to convert each patient into a vector of pathway
activity scores~\cite{gillespie2022reactome}. This representation is easier
to connect back to biology because each feature corresponds to a known
pathway rather than an isolated gene.

We formulate our problem as a multi-label prediction task. Given one
patient's Reactome pathway activity profile, the model predicts three binary
outcomes: targeted molecular therapy (TMT), radiation therapy (RT), and
overall survival of at least 180 days (OS $\geq$ 180 d). This setup is
important because a patient can belong to any combination of these outcomes.
For example, one patient may receive targeted therapy but not radiation,
while another may receive radiation and also survive beyond 180 days. By
using three output heads, the model can learn which pathway patterns are
most useful for each clinical question.

Recent work has shown that biological knowledge can be built directly into
deep learning models. P-NET used pathway structure to create a neural
network for prostate cancer discovery~\cite{elmarakeby2021pnet}. BINN
extended this idea by using biologically informed network connections to
support pathway-level interpretation~\cite{hartman2023binn}. GraphPath
models pathways as nodes in a graph and uses graph attention to learn
relationships between pathways~\cite{ma2024graphpath,velickovic2018gat}.
PATH uses a graph transformer to represent pathway interactions for cancer
prognosis~\cite{howlader2026path}. These models are promising, but they
were developed and evaluated under different settings. Their datasets,
labels, input features, and evaluation tasks are not the same, so it is
difficult to know whether one architecture is generally better or whether
each one works best for different clinical outcomes.

In this work, we compare these pathway-informed models under one common
pipeline. We use gene expression and clinical data from The Cancer Genome
Atlas (TCGA), accessed through the UCSC Xena platform~\cite{Weinstein2013,
xena2020}. We study five solid-tumour cohorts: breast, lung, prostate, head
and neck, and thyroid cancer. For every cohort, we use the same preprocessing
strategy: gene expression is converted into Reactome pathway activity scores,
clinical records are converted into the three binary labels, and the same
multi-label learning objective is used for all models. This makes the
comparison more direct because the main difference between models is their
architecture, not the data pipeline.

Our goal is not only to report which model has the highest score, but also
to understand when each kind of biological structure helps. BINN represents
Reactome as a sparse hierarchy. GraphPath represents pathways as an
attention-based graph. PATH uses a transformer-style graph model with
pathway interactions. By testing these approaches on the same task, we ask a
simple question: do different pathway-based architectures help with
different therapy and survival outcomes? This question matters because label
distributions are different across cancer types. A cohort with many patients
receiving targeted therapy may not behave like a cohort where targeted
therapy is rare. Therefore, the best model may depend on both the phenotype
being predicted and the cancer cohort being studied.

The main contributions of this paper are as follows:
\begin{itemize}[leftmargin=*]
  \item We build a unified TCGA benchmark for pathway-based prediction of
  TMT, RT, and OS $\geq$ 180 d across five solid-tumour cohorts.
  \item We use the same Reactome and ssGSEA-based pathway representation for
  every patient, making the input space consistent across cohorts and models.
  \item We adapt and compare three pathway-informed deep learning
  architectures: BINN, GraphPath, and PATH.
  \item We show that model performance is phenotype-specific, meaning that
  the strongest architecture can change depending on whether the target is
  therapy assignment or short-term survival.
\end{itemize}
